\section{Drift Detection Results}
\label{sec:5_6}

DDM and ADWIN monitors tracked prediction quality across the evaluation period.

\textbf{DDM Findings:}

\begin{table}[htbp]
    \centering
    \caption{DDM Drift Detection}
    \label{tab:5_7}
    \begin{tabular}{lccc}
        \toprule
        Period & Warning Triggered & Drift Detected & Error Rate \\
        \midrule
        Feb 2025 & Yes & No & 8.2\% \\
        Mar 2025 & Yes & Yes & 11.4\% \\
        Apr 2025 & No & No & 6.8\% \\
        May 2025 & No & No & 5.9\% \\
        \bottomrule
    \end{tabular}
\end{table}

DDM detected significant drift in March 2025, coinciding with the fraud rate spike identified in EDA. The 3$\sigma$ threshold triggered retraining alerts.

\textbf{ADWIN Findings:}

ADWIN's adaptive windowing identified distribution changes in:
\begin{itemize}
    \item March 2025: Fraud pattern shift (confirmed by DDM)
    \item June 2025: Gradual drift in feature distributions
\end{itemize}

ADWIN's parameter-free approach detected gradual drift earlier than DDM's fixed thresholds.

\textbf{Sliding vs Expanding Window:}

\begin{table}[htbp]
    \centering
    \caption{Window Strategy Comparison}
    \label{tab:5_8}
    \begin{tabular}{lccc}
        \toprule
        Validation & Mean AUC-PR & Std & Drift Range \\
        \midrule
        Expanding (accumulating) & 0.6286 & 0.077 & 0.21 \\
        Sliding (90-day fixed) & 0.5892 & 0.092 & 0.28 \\
        Sliding (180-day fixed) & 0.6104 & 0.081 & 0.23 \\
        \bottomrule
    \end{tabular}
\end{table}

Expanding window outperforms sliding windows, indicating historical data provides value despite potential staleness. The 90-day sliding window shows highest variance, insufficient for stable model performance.

\textbf{Recommendation:} Use expanding window training with DDM/ADWIN monitoring for retraining triggers. Monthly scheduled retraining provides good balance between freshness and stability.
